% GENERATED FILE. Edit canonical lessons, then run npm run generate:books.
% canonical-source-hash: fnv1a64:465ab8b17e46fdd9
% canonical-lessons: KA-C532-gite, KA-C532-nrtya, KA-C532-hasya, KA-C532-jogula, KA-C532-meluku-haku, KA-R532-review-526-532

\chapter{Song, Dance and Laughter}
\label{ch:ka-song-dance-and-laughter}
\hlchaptermodality{532}

\begin{chapteropening}
\textbf{By the end of this chapter:} \emph{I can say a song, a dance form, humour, a lullaby and to reminisce.}

Complete the last lesson of chapter 532: I can say a song, a dance form, humour, a lullaby and to reminisce.
\end{chapteropening}

\section[\texorpdfstring{\kn{ಗೀತೆ} (gīte)}{ಗೀತೆ (gīte)}]{\kn{ಗೀತೆ} (gīte) — a song}

\label{lesson:KA-C532-gite}



\begin{quote}
Before the new one: say the Kannada for a swing, then the Kannada for a toy.
\end{quote}

\subsection*{You'll want to know: \kn{ಗೀತೆ}}
\textbf{\kn{ಗೀತೆ}} — \emph{gīte} — \textquotedblleft{}a song\textquotedblright{}.

\begin{quote}
\kn{ಶಾಲೆಯ} \kn{ಸಮಾರಂಭದಲ್ಲಿ} \kn{ಮಕ್ಕಳು} \kn{ಒಂದು} \kn{ಗೀತೆ} \kn{ಹಾಡಿದರು}.
\end{quote}

\begin{quote}
\emph{śāleya samārambhadalli makkaḷu ondu gīte hāḍidaru.}
\end{quote}

\begin{quote}
The children sang a song at the school function.
\end{quote}

Formal; in talk people say \kn{ಹಾಡು}. \kn{ನಾಡಗೀತೆ}, the state song.

\begin{cousinweb}
\textbf{\kn{ಗೀತೆ}} (\emph{gīte}) goes back to Sanskrit \emph{gīta}, \textquotedblleft{}singing\textquotedblright{} (CDIAL 4167), from \emph{gai}, \textquotedblleft{}to sing\textquotedblright{}.
\end{cousinweb}

\begin{grammarlens}[title={What you've built}]
One more word for describing what you lived and felt.
\end{grammarlens}

\subsection*{Your turn}
\begin{itemize}
\raggedright
  \item \emph{Say it:} \emph{gīte}
  \item \emph{Say it:} \emph{gīte}, once more
  \item \emph{Say it:} \emph{śāleya samārambhadalli makkaḷu ondu gīte hāḍidaru.} — The children sang a song at the school function.
  \item \emph{Recall:} say the Kannada for a swing, then the Kannada for a toy, then say \emph{gīte} again
\end{itemize}

\begin{tcolorbox}[breakable,colback=teal!4,colframe=teal!35!black,title={Before you move on}]
What does \kn{ಗೀತೆ} mean? (\textquotedblleft{}A song\textquotedblright{}.) Say it once more.
\end{tcolorbox}
\section[\texorpdfstring{\kn{ನೃತ್ಯ} (nṛtya)}{ನೃತ್ಯ (nṛtya)}]{\kn{ನೃತ್ಯ} (nṛtya) — a dance form; dancing}

\label{lesson:KA-C532-nrtya}



\begin{quote}
Before the new one: say the Kannada for a toy, then the Kannada for a song.
\end{quote}

\subsection*{You'll want to know: \kn{ನೃತ್ಯ}}
\textbf{\kn{ನೃತ್ಯ}} — \emph{nṛtya} — \textquotedblleft{}a dance form; dancing\textquotedblright{}.

\begin{quote}
\kn{ಅಕ್ಕ} \kn{ಹತ್ತು} \kn{ವರ್ಷದಿಂದ} \kn{ನೃತ್ಯ} \kn{ಕಲಿಯುತ್ತಿದ್ದಾಳೆ}.
\end{quote}

\begin{quote}
\emph{akka hattu varṣadinda nṛtya kaliyuttiddāḷe.}
\end{quote}

\begin{quote}
My elder sister has been learning dance for ten years.
\end{quote}

Formal and artistic; everyday dancing is \kn{ಕುಣಿತ}.

\begin{cousinweb}
\textbf{\kn{ನೃತ್ಯ}} (\emph{nṛtya}) is borrowed from Sanskrit \emph{nṛtya}, \textquotedblleft{}dancing, acting, gesture\textquotedblright{} (Monier-Williams).
\end{cousinweb}

\begin{grammarlens}[title={What you've built}]
One more word for describing what you lived and felt.
\end{grammarlens}

\subsection*{Your turn}
\begin{itemize}
\raggedright
  \item \emph{Say it:} \emph{nṛtya}
  \item \emph{Say it:} \emph{nṛtya}, once more
  \item \emph{Say it:} \emph{akka hattu varṣadinda nṛtya kaliyuttiddāḷe.} — My elder sister has been learning dance for ten years.
  \item \emph{Recall:} say the Kannada for a toy, then the Kannada for a song, then say \emph{nṛtya} again
\end{itemize}

\begin{tcolorbox}[breakable,colback=teal!4,colframe=teal!35!black,title={Before you move on}]
What does \kn{ನೃತ್ಯ} mean? (\textquotedblleft{}A dance form; dancing\textquotedblright{}.) Say it once more.
\end{tcolorbox}
\section[\texorpdfstring{\kn{ಹಾಸ್ಯ} (hāsya)}{ಹಾಸ್ಯ (hāsya)}]{\kn{ಹಾಸ್ಯ} (hāsya) — humour; comedy}

\label{lesson:KA-C532-hasya}



\begin{quote}
Before the new one: say the Kannada for a song, then the Kannada for a dance form.
\end{quote}

\subsection*{You'll want to know: \kn{ಹಾಸ್ಯ}}
\textbf{\kn{ಹಾಸ್ಯ}} — \emph{hāsya} — \textquotedblleft{}humour; comedy\textquotedblright{}.

\begin{quote}
\kn{ಆ} \kn{ಸಿನಿಮಾದಲ್ಲಿ} \kn{ತುಂಬಾ} \kn{ಹಾಸ್ಯ} \kn{ಇದೆ}.
\end{quote}

\begin{quote}
\emph{ā sinimādalli tumbā hāsya ide.}
\end{quote}

\begin{quote}
There is a lot of humour in that film.
\end{quote}

\kn{ಹಾಸ್ಯಗಾರ}, a comedian.

\begin{cousinweb}
\textbf{\kn{ಹಾಸ್ಯ}} (\emph{hāsya}) is borrowed from Sanskrit \emph{hāsya}, \textquotedblleft{}laughable, funny; laughter\textquotedblright{} (CDIAL 14071).
\end{cousinweb}

\begin{grammarlens}[title={What you've built}]
One more word for describing what you lived and felt.
\end{grammarlens}

\subsection*{Your turn}
\begin{itemize}
\raggedright
  \item \emph{Say it:} \emph{hāsya}
  \item \emph{Say it:} \emph{hāsya}, once more
  \item \emph{Say it:} \emph{ā sinimādalli tumbā hāsya ide.} — There is a lot of humour in that film.
  \item \emph{Recall:} say the Kannada for a song, then the Kannada for a dance form, then say \emph{hāsya} again
\end{itemize}

\begin{tcolorbox}[breakable,colback=teal!4,colframe=teal!35!black,title={Before you move on}]
What does \kn{ಹಾಸ್ಯ} mean? (\textquotedblleft{}Humour; comedy\textquotedblright{}.) Say it once more.
\end{tcolorbox}
\section[\texorpdfstring{\kn{ಜೋಗುಳ} (jōguḷa)}{ಜೋಗುಳ (jōguḷa)}]{\kn{ಜೋಗುಳ} (jōguḷa) — a lullaby}

\label{lesson:KA-C532-jogula}



\begin{quote}
Before the new one: say the Kannada for a dance form, then the Kannada for humour.
\end{quote}

\subsection*{You'll want to know: \kn{ಜೋಗುಳ}}
\textbf{\kn{ಜೋಗುಳ}} — \emph{jōguḷa} — \textquotedblleft{}a lullaby\textquotedblright{}.

\begin{quote}
\kn{ಅಮ್ಮ} \kn{ಜೋಗುಳ} \kn{ಹಾಡಿದಾಗ} \kn{ಮಗು} \kn{ಮಲಗಿತು}.
\end{quote}

\begin{quote}
\emph{amma jōguḷa hāḍidāga magu malagitu.}
\end{quote}

\begin{quote}
When the mother sang a lullaby, the baby fell asleep.
\end{quote}

Built on \kn{ಜೋ} \kn{ಜೋ}, the soothing sound sung over a cradle.

\begin{cousinweb}
\textbf{\kn{ಜೋಗುಳ}} (\emph{jōguḷa}) is native Dravidian, built on Kannada \emph{jō}, \textquotedblleft{}hush!\textquotedblright{}, the word sung in lullabies (DEDR 2869).
\end{cousinweb}

\begin{grammarlens}[title={What you've built}]
One more word for describing what you lived and felt.
\end{grammarlens}

\subsection*{Your turn}
\begin{itemize}
\raggedright
  \item \emph{Say it:} \emph{jōguḷa}
  \item \emph{Say it:} \emph{jōguḷa}, once more
  \item \emph{Say it:} \emph{amma jōguḷa hāḍidāga magu malagitu.} — When the mother sang a lullaby, the baby fell asleep.
  \item \emph{Recall:} say the Kannada for a dance form, then the Kannada for humour, then say \emph{jōguḷa} again
\end{itemize}

\begin{tcolorbox}[breakable,colback=teal!4,colframe=teal!35!black,title={Before you move on}]
What does \kn{ಜೋಗುಳ} mean? (\textquotedblleft{}A lullaby\textquotedblright{}.) Say it once more.
\end{tcolorbox}
\section[\texorpdfstring{\kn{ಮೆಲುಕು} \kn{ಹಾಕು} (meluku hāku)}{ಮೆಲುಕು ಹಾಕು (meluku hāku)}]{\kn{ಮೆಲುಕು} \kn{ಹಾಕು} (meluku hāku) — to reminisce; to chew over}

\label{lesson:KA-C532-meluku-haku}



\begin{quote}
Before the new one: say the Kannada for humour, then the Kannada for a lullaby.
\end{quote}

\subsection*{You'll want to know: \kn{ಮೆಲುಕು} \kn{ಹಾಕು}}
\textbf{\kn{ಮೆಲುಕು} \kn{ಹಾಕು}} — \emph{meluku hāku} — \textquotedblleft{}to reminisce; to chew over\textquotedblright{}.

\begin{quote}
\kn{ಅಜ್ಜ} \kn{ತಮ್ಮ} \kn{ಬಾಲ್ಯದ} \kn{ದಿನಗಳನ್ನು} \kn{ಮೆಲುಕು} \kn{ಹಾಕಿದರು}.
\end{quote}

\begin{quote}
\emph{ajja tamma bālyada dinagaḷannu meluku hākidaru.}
\end{quote}

\begin{quote}
Grandfather went back over his childhood days.
\end{quote}

Past (I …): \kn{ಮೆಲುಕು} \kn{ಹಾಕಿದೆ}, \emph{meluku hākide}.

\kn{ಮೆಲುಕು} is a cow chewing the cud: going back over memories again and again.

\begin{cousinweb}
\textbf{\kn{ಮೆಲುಕು} \kn{ಹಾಕು}} (\emph{meluku hāku}) is a native Kannada phrase: \emph{meluku}, \textquotedblleft{}bringing food up again to chew the cud\textquotedblright{} (DEDR 5077), with \kn{ಹಾಕು}, \textquotedblleft{}to put\textquotedblright{}.
\end{cousinweb}

\begin{grammarlens}[title={What you've built}]
One more word for describing what you lived and felt.
\end{grammarlens}

\subsection*{Your turn}
\begin{itemize}
\raggedright
  \item \emph{Say it:} \emph{meluku hāku}
  \item \emph{Say it:} \emph{meluku hāku}, once more
  \item \emph{Say it:} \emph{ajja tamma bālyada dinagaḷannu meluku hākidaru.} — Grandfather went back over his childhood days.
  \item \emph{Recall:} say the Kannada for humour, then the Kannada for a lullaby, then say \emph{meluku hāku} again
\end{itemize}

\begin{tcolorbox}[breakable,colback=teal!4,colframe=teal!35!black,title={Before you move on}]
What does \kn{ಮೆಲುಕು} \kn{ಹಾಕು} mean? (\textquotedblleft{}To reminisce; to chew over\textquotedblright{}.) Say it once more.
\end{tcolorbox}
\section[(dialogue)]{Looking back: from Kin to Song, Dance and Laughter}

\label{lesson:KA-R532-review-526-532}



\begin{quote}
Say the Kannada for kin and a woman's parental home.
\end{quote}

\subsection*{Your turn}
Say these aloud:

\begin{itemize}
\raggedright
  \item kin, a woman's parental home, twins, a great-grandson, younger
  \item to adopt, to bring up, to nourish, crawling on all fours, a stepmother
  \item friendship, a bond, companionship, dear, to make up
  \item an adult, the teenage years, solitude, infatuation, to live
  \item a pause, entertainment, a chat, to savour, to be absorbed
  \item a dice game, carrom, hide-and-seek, a swing, a toy
  \item a song, a dance form, humour, a lullaby, to reminisce
\end{itemize}

\begin{tcolorbox}[breakable,colback=teal!4,colframe=teal!35!black,title={Before you move on}]
Kin? (\textbf{\kn{ಬಂಧು}}, \emph{bandhu}.) A woman's parental home? (\textbf{\kn{ತವರು}}, \emph{tavaru}.) Twins? (\textbf{\kn{ಅವಳಿ}}, \emph{avaḷi}.) A great-grandson? (\textbf{\kn{ಮರಿಮಗ}}, \emph{marimaga}.) Younger? (\textbf{\kn{ಕಿರಿಯ}}, \emph{kiriya}.) To adopt? (\textbf{\kn{ದತ್ತು} \kn{ತೆಗೆದುಕೊ}}, \emph{dattu tegeduko}.) To bring up? (\textbf{\kn{ಸಲಹು}}, \emph{salahu}.) To nourish? (\textbf{\kn{ಪೋಷಿಸು}}, \emph{pōṣisu}.) Crawling on all fours? (\textbf{\kn{ಅಂಬೆಗಾಲು}}, \emph{ambegālu}.) A stepmother? (\textbf{\kn{ಮಲತಾಯಿ}}, \emph{malatāyi}.) Friendship? (\textbf{\kn{ಸ್ನೇಹ}}, \emph{snēha}.) A bond? (\textbf{\kn{ಬಾಂಧವ್ಯ}}, \emph{bāndhavya}.) Companionship? (\textbf{\kn{ಒಡನಾಟ}}, \emph{oḍanāṭa}.) Dear? (\textbf{\kn{ಆತ್ಮೀಯ}}, \emph{ātmīya}.) To make up? (\textbf{\kn{ರಾಜಿ} \kn{ಮಾಡಿಕೊ}}, \emph{rāji māḍiko}.) An adult? (\textbf{\kn{ವಯಸ್ಕ}}, \emph{vayaska}.) The teenage years? (\textbf{\kn{ಹದಿಹರೆಯ}}, \emph{hadihareya}.) Solitude? (\textbf{\kn{ಏಕಾಂತ}}, \emph{ēkānta}.) Infatuation? (\textbf{\kn{ಮೋಹ}}, \emph{mōha}.) To live? (\textbf{\kn{ಬದುಕು}}, \emph{baduku}.) A pause? (\textbf{\kn{ವಿರಾಮ}}, \emph{virāma}.) Entertainment? (\textbf{\kn{ಮನರಂಜನೆ}}, \emph{manarañjane}.) A chat? (\textbf{\kn{ಹರಟೆ}}, \emph{haraṭe}.) To savour? (\textbf{\kn{ಆಸ್ವಾದಿಸು}}, \emph{āsvādisu}.) To be absorbed? (\textbf{\kn{ಮೈಮರೆ}}, \emph{maimare}.) A dice game? (\textbf{\kn{ಪಗಡೆ}}, \emph{pagaḍe}.) Carrom? (\textbf{\kn{ಕೇರಂ}}, \emph{kēraṃ}.) Hide-and-seek? (\textbf{\kn{ಕಣ್ಣಾಮುಚ್ಚಾಲೆ}}, \emph{kaṇṇāmuccāle}.) A swing? (\textbf{\kn{ಉಯ್ಯಾಲೆ}}, \emph{uyyāle}.) A toy? (\textbf{\kn{ಆಟಿಕೆ}}, \emph{āṭike}.) A song? (\textbf{\kn{ಗೀತೆ}}, \emph{gīte}.) A dance form? (\textbf{\kn{ನೃತ್ಯ}}, \emph{nṛtya}.) Humour? (\textbf{\kn{ಹಾಸ್ಯ}}, \emph{hāsya}.) A lullaby? (\textbf{\kn{ಜೋಗುಳ}}, \emph{jōguḷa}.) To reminisce? (\textbf{\kn{ಮೆಲುಕು} \kn{ಹಾಕು}}, \emph{meluku hāku}.)
\end{tcolorbox}
